\Needspace{16\baselineskip}
\section{Input feature order}
\label{app:features}
Table~\ref{tab:features} defines the coordinates $r_f$ of the archived feature
vector, using the source dataset's feature names. All displayed cuts use their
standardized counterparts $x_f$ from Eq.~\eqref{eq:standardization}.
\begin{table}[H]\centering\small
\caption{Feature-coordinate convention. Names identify columns in the
high-level jet dataset~\cite{openml42468}.}\label{tab:features}
\begin{tabular}{@{}rlrl@{}}\toprule
$f$ & Dataset feature & $f$ & Dataset feature\\\midrule
1 & \texttt{zlogz} & 9 & \texttt{d2\_a1\_b1\_mmdt}\\
2 & \texttt{c1\_b0\_mmdt} & 10 & \texttt{d2\_a1\_b2\_mmdt}\\
3 & \texttt{c1\_b1\_mmdt} & 11 & \texttt{m2\_b1\_mmdt}\\
4 & \texttt{c1\_b2\_mmdt} & 12 & \texttt{m2\_b2\_mmdt}\\
5 & \texttt{c2\_b1\_mmdt} & 13 & \texttt{n2\_b1\_mmdt}\\
6 & \texttt{c2\_b2\_mmdt} & 14 & \texttt{n2\_b2\_mmdt}\\
7 & \texttt{d2\_b1\_mmdt} & 15 & \texttt{mass\_mmdt}\\
8 & \texttt{d2\_b2\_mmdt} & 16 & \texttt{multiplicity}\\\bottomrule
\end{tabular}
\end{table}
